% Vista editorial exacta de corpus/source/masas/06_operadores.tex, líneas 50--113. Las líneas 1--29 y 30--49 ya residen exactamente en 73_rev2_electron_lectores_torres.tex.
\subsection{Lecteurs bidirectionnels \(K_D\) et \(K_\Omega\) : profils, coïncidences et opérateurs par multisection}
Le parcours orienté produit deux familles d'invariants. Pour les défauts temporels \(D_j\) et la mémoire de retour \(\Omega\), nous définissons

\[
K_D=\left(D_{27}+D_{36},D_1,\frac{D_4-D_3}{2}\right),
\qquad
K_\Omega=(\Omega,54,P_6).
\]
Sur les 324 parcours apparaissent 14 profils distincts de \(K_D\), 9 profils de \(K_\Omega\), 40 paires conjointes et 18 coïncidences. L'unique profil commun du parcours électronique est

\[
K_D(\Gamma_e)=K_\Omega(\Gamma_e)=(80,54,6).
\]
Ces lecteurs ne sont pas des étiquettes décoratives. \(K_D\) mesure les écarts du mot ; \(K_\Omega\) mesure le retour enrichi. Leur égalité au centre certifie la clôture bidirectionnelle de rang un. Hors du centre, ils produisent des opérateurs de fibre dont les spectres complets doivent être présentés, même lorsqu'une réalisation physique ultérieure sélectionne une section scalaire.

\Needspace{7\baselineskip}
\subsection{Raffinement harmonique propre à chaque état}
Le raffinement complet se situe dans l'espace de coefficients

\[
\mathcal E_{90,120}
=\left\{\eta=(\eta_h)_{h\in\langle90,120\rangle}:
\sum_h|\eta_hs_h|<\infty\right\}.
\]
L'application recherchée ne reçoit pas une masse. Elle reçoit le parcours enrichi :

\[
\mathfrak T:\Gamma_{\rm enr}\longrightarrow\mathcal E_{90,120},
\qquad
\Gamma\longmapsto\eta(\Gamma).
\]
Ses coordonnées doivent être des fonctions explicites des retours, de la retenue, de l'enroulement, de la mémoire et de la frontière. La condition de naturalité est

\[
\eta(r_{n,m}\Gamma)=r^{\mathcal E}_{n,m}\eta(\Gamma),
\]
et la convergence doit s'accompagner d'une borne de reste calculable. Le caractère complet agit alors par

\[
\widehat M_F
=\sum_{\gamma\in F}
m_{\rm clk}\,
\chi_{\rm full}(\sigma(\gamma),\eta(\gamma))
|\gamma\rangle\langle\gamma|.
\]
La capacité de reconstruire une différence numérique donnée démontre la complétude représentative de la hiérarchie ; elle ne démontre pas, à elle seule, que \(\mathfrak T\) a choisi ses coefficients. La différence entre les deux affirmations est préservée dans chaque ligne quantitative.

\Needspace{7\baselineskip}
\subsection{Scalarisation par couplage}
Soit \(P_{X,a}\) le projecteur déterminé par l'état \(X\) et l'appareil ou le canal physique \(a\). La distribution de masse est

\[
\mu_{X,a}(B)
=\langle\Psi_X,
P_{X,a}E_{\widehat M_F}(B)P_{X,a}\Psi_X\rangle.
\]
Il existe une droite scalaire lorsque la compression satisfait

\[
P_{X,a}\widehat M_FP_{X,a}=m_XP_{X,a}.
\]
Tel est le sens précis de « mesurer, c'est se coupler ». La valeur numérique n'était pas cachée dans une cellule nominale ; elle apparaît comme valeur propre d'une compression physique d'un opérateur déjà construit. Si la compression conserve plusieurs valeurs propres, la prédiction correcte est un multiplet ou une distribution, et non une moyenne choisie pour coïncider avec une référence.

\Needspace{7\baselineskip}

